% !PRINTED 14
\weberpage{14}

%% [running head, not body text -- commented out 2026-08-07] Page furniture (running
%% head, folio, signature) is excluded from this edition by standing policy.
%% \gesperrt{Erster Abschnitt.} \hfill 14 \hfill \gesperrt{\S.\,4.}

Wegen der Eigenschaft der Normaltheiler ist aber $N$ mit jedem $b$ vertauschbar, also $b_h N = N b_h$, und folglich
\[
N_h N_k = N N b_h b_k.
\]

Weil $b_h b_k$ in $P$ enthalten ist, so ist $N b_h b_k$ mit einer der Nebengruppen (5) identisch, und da $NN = N$ gesetzt werden kann (nach 1.), so folgt, wenn wir $N b_h b_k = N_l$ setzen,
\weq{6}{N_h N_k = N_l.}

Damit ist die Eigenschaft \S.\,1, 1. der Gruppen für die Composition der Nebengruppen nachgewiesen. Dass auch \S.\,1, 2., nämlich das associative Gesetz gilt, haben wir schon hervorgehoben. Es ist also noch \S.\,1, 3. nachzuweisen, nämlich dass, wenn
\weq{7}{N_h N_i = N_k N_l}
gesetzt wird, und $N_h = N_k$ ist, auch $N_i = N_l$, und wenn $N_i = N_l$ ist, auch $N_h = N_k$ sein muss.

Nach der Darstellung (5) können wir aber (7) in der doppelten Weise darstellen:
\weq{8}{\begin{aligned}
N b_h b_i &= N b_k b_l \\
b_h b_i N &= b_k b_l N.
\end{aligned}}

Ist nun $N_i = N_l$, so können wir auch $b_i = b_l$ annehmen, und erhalten aus (8) durch Composition mit $b_i^{-1}$
\[
N b_h = N b_k;
\]
und wenn $N b_h = N b_k$ ist, so nehmen wir $b_h = b_k$ an und erhalten aus (8) durch Composition mit $b_h^{-1}$
\[
b_i N = b_l N,
\]
also auch $N_i = N_l$.

Damit ist also erwiesen, dass das System der Nebengruppen durch die Composition der Theile zu einer Gruppe wird; dass in dieser Gruppe der Normaltheiler $N$ das Einheitselement ist, folgt unmittelbar aus 1., weil danach
\[
N N_k = N N b_k = N_k
\]
ist. Und weil $N b_k N b_k^{-1} = N b_k b_k^{-1} = N$ ist, so sind $N b_k$ und $N b_k^{-1}$ entgegengesetzte Elemente.

Die Gruppe der Grössen $N_k$, deren Existenz also hiermit nachgewiesen ist, nennen wir die \gesperrt{Gruppe der Nebengruppen} oder die \gesperrt{zu $N$ complementäre Gruppe} in Bezug auf $P$. Wir
%CONT>>
